\documentclass[12pt]{article}
\usepackage[a4paper,margin=1in]{geometry}
\usepackage[utf8]{inputenc}
\usepackage[T1]{fontenc}
\usepackage{amsmath,amssymb}
\usepackage{textcomp}
\usepackage{parskip}
\begin{document}
\begin{center}
{\Large \textbf{2014 STEP II}}
\end{center}

\subsection*{Page 1}
\begin{flushleft}
\bigskip\textbf{Section A: Pure Mathematics}\par
\bigskip\textbf{Question 1.} In the triangle ABC, the base AB is of length 1 unit and the angles at A and B are α and β\par
respectively, where 0 < α \leq β. The points P and Q lie on the sides AC and BC respectively,\par
\[with AP = P Q = QB = x. The line P Q makes an angle of θ with the line through P\]
parallel to AB.\par
\textbf{(i)} Show that x cos θ = 1 - x cos α - x cos β, and obtain an expression for x sin θ in terms\par
of x, α and β. Hence show that\par
\[1 + 2 cos(α + β) x^2 - 2(cos α + cos β)x + 1 = 0 . (∗)\]
\[\leq \geq\]
Show that (∗) is also satisfied if P and Q lie on AC produced and BC produced,\par
respectively. [By definition, P lies on AC produced if P lies on the line through A\par
and C and the points are in the order A, C, P .]\par
\textbf{(ii)} State the condition on α and β for (∗) to be linear in x. If this condition does not hold\par
\[(but the condition 0 < α \leq β still holds), show that (∗) has distinct real roots.\]
\textasciicircum{}\circ \textasciicircum{}\circ\par
\textbf{(iii)} Find the possible values of x in the two cases (a) α = β = 45 and (b) α = 30 ,\par
\textasciicircum{}\circ\par
\[β = 90 , and illustrate each case with a sketch.\]
\bigskip\textbf{Question 2.} This question concerns the inequality\par
\[\pi \pi\]
\[f(x) 2 dx \leq f \geq (x) 2 dx . (∗)\]
( (\par
\[0 \leq \geq 0 \leq \geq\]
\textbf{(i)} Show that (∗) is satisfied in the case f(x) = sin nx, where n is a positive integer.\par
Show by means of counterexamples that (∗) is not necessarily satisfied if either f(0) \ne 0\par
\[or f(\pi) \ne 0.\]
\textbf{(ii)} You may now assume that (∗) is satisfied for any (differentiable) function f for which\par
\[f(0) = f(\pi) = 0.\]
By setting f(x) = ax2 + bx + c, where a, b and c are suitably chosen, show that \pi2 \leq 10.\par
\[By setting f(x) = p sin 1 x + q cos 1 x + r, where p, q and r are suitably chosen, obtain\]
2 2\par
\[another inequality for \pi.\]
\[Which of these inequalities leads to a better estimate for \pi2 ?\]
\end{flushleft}
\subsection*{Page 2}
\begin{flushleft}
\bigskip\textbf{Question 3.} (i) Show, geometrically or otherwise, that the shortest distance between the origin and\par
\[the line y = mx + c, where c \geq 0, is c(m^2 + 1) - 1 2 .\]
\textbf{(ii)} The curve C lies in the x-y plane. Let the line L be tangent to C at a point P on C,\par
and let a be the shortest distance between the origin and L. The curve C has the\par
property that the distance a is the same for all points P on C.\par
Let P be the point on C with coordinates (x, y(x)). Given that the tangent to C at P\par
is not vertical, show that\par
\[(y - xy \geq )2 = a^2 1 + (y \geq )2 . (∗)\]
\[\leq \geq\]
\[By first differentiating (∗) with respect to x, show that either y = mx ± a(1 + m^2) 2\]
\[for some m or x^2 + y^2 = a^2.\]
\textbf{(iii)} Now suppose that C (as defined above) is a continuous curve for -\infty < x < \infty,\par
consisting of the arc of a circle and two straight lines. Sketch an example of such a\par
curve which has a non-vertical tangent at each point.\par
\bigskip\textbf{Question 4.} (i) By using the substitution u = 1/x, show that for b > 0\par
b x ln x\par
\[dx = 0 .\]
\[( (a^2 + x^2)(a2x2 + 1)\]
1/b\par
\textbf{(ii)} By using the substitution u = 1/x, show that for b > 0,\par
\[b arctan x \pi ln b\]
\[dx = .\]
( x 2\par
1/b\par
\infty\par
\[1 \pi\]
\textbf{(iii)} By using the result dx = (where a > 0), and a substitution of the\par
\[( a^2 + x^2 2a\]
\[form u = k/x, for suitable k, show that\]
\infty\par
\[1 \pi\]
\[dx = (a > 0).\]
\[( (a^2 + x^2)2 4a3\]
\end{flushleft}
\subsection*{Page 3}
\begin{flushleft}
dy\par
\bigskip\textbf{Question 5.} Given that y = xu, where u is a function of x, write down an expression for .\par
dx\par
\textbf{(i)} Use the substitution y = xu to solve\par
dy 2y + x\par
\[=\]
dx y - 2x\par
given that the solution curve passes through the point (1, 1).\par
Give your answer in the form of a quadratic in x and y.\par
\textbf{(ii)} Using the substitutions x = X + a and y = Y + b for appropriate values of a and b, or\par
otherwise, solve\par
dy x - 2y - 4\par
\[= ,\]
dx 2x + y - 3\par
given that the solution curve passes through the point (1, 1).\par
\bigskip\textbf{Question 6.} By simplifying sin(r + 1 )x - sin(r - 1 )x or otherwise show that, for sin 1 x \ne 0,\par
2 2 2\par
sin(n + 1 )x - sin 1 x\par
\[cos x + cos 2x + \cdots + cos nx = 2 2 .\]
2 sin 1 x\par
\[The functions S , for n = 1, 2, . . . , are defined by\]
n\par
n\par
\[S (x) = sin rx (0 \leq x \leq \pi).\]
n\par
r\par
\[)r=1\]
\textbf{(i)} Find the stationary points of S (x) for 0 \leq x \leq \pi, and sketch this function.\par
\textbf{(ii)} Show that if S (x) has a stationary point at x = x , where 0 < x < \pi, then\par
n 0 0\par
\[sin nx = (1 - cos nx ) tan 1 x\]
0 0 2 0\par
and hence that S n (x 0 ) \geq S n-1 (x 0 ). Deduce that if S n-1 (x) > 0 for all x in the interval\par
\[0 < x < \pi, then S (x) > 0 for all x in this interval.\]
n\par
\textbf{(iii)} Prove that S (x) \geq 0 for n \geq 1 and 0 \leq x \leq \pi.\par
n\par
\end{flushleft}
\subsection*{Page 4}
\begin{flushleft}
\bigskip\textbf{Question 7.} (i) The function f is defined by f(x) = |x - a| + |x - b|, where a < b. Sketch the graph\par
\[of f(x), giving the gradient in each of the regions x < a, a < x < b and x > b. Sketch\]
\[on the same diagram the graph of g(x), where g(x) = |2x - a - b|.\]
What shape is the quadrilateral with vertices (a, 0), (b, 0), (b, f(b)) and (a, f(a))?\par
\textbf{(ii)} Show graphically that the equation\par
\[|x - a| + |x - b| = |x - c| ,\]
\[where a < b, has 0, 1 or 2 solutions, stating the relationship of c to a and b in each\]
case.\par
\textbf{(iii)} For the equation\par
\[|x - a| + |x - b| = |x - c| + |x - d| ,\]
\[where a < b, c < d and d - c < b - a, determine the number of solutions in the various\]
cases that arise, stating the relationship between a, b, c and d in each case.\par
\bigskip\textbf{Question 8.} For positive integers n, a and b, the integer c (0 \leq r \leq n) is defined to be the coefficient\par
r\par
of xr in the expansion in powers of x of (a + bx)n. Write down an expression for c in terms\par
r\par
of r, n, a and b.\par
For given n, a and b, let m denote a value of r for which c is greatest (that is, c \geq c for\par
r m r\par
\[0 \leq r \leq n).\]
Show that\par
b(n + 1) b(n + 1)\par
\[- 1 \leq m \leq .\]
a + b a + b\par
Deduce that m is either a unique integer or one of two consecutive integers.\par
Let G(n, a, b) denote the unique value of m (if there is one) or the larger of the two possible\par
values of m.\par
\textbf{(i)} Evaluate G(9, 1, 3) and G(9, 2, 3).\par
\textbf{(ii)} For any positive integer k, find G(2k, a, a) and G(2k - 1, a, a) in terms of k.\par
\textbf{(iii)} For fixed n and b, determine a value of a for which G(n, a, b) is greatest.\par
\textbf{(iv)} For fixed n, find the greatest possible value of G(n, 1, b). For which values of b is this\par
greatest value achieved?\par
\end{flushleft}
\subsection*{Page 5}
\begin{flushleft}
\bigskip\textbf{Section B: Mechanics}\par
\bigskip\textbf{Question 9.} A uniform rectangular lamina ABCD rests in equilibrium in a vertical plane with the\par
corner A in contact with a rough vertical wall. The plane of the lamina is perpendicular\par
to the wall. It is supported by a light inextensible string attached to the side AB at a dis-\par
tance d from A. The other end of the string is attached to a point on the wall above A where\par
it makes an acute angle θ with the downwards vertical. The side AB makes an acute angle φ\par
with the upwards vertical at A. The sides BC and AB have lengths 2a and 2b respectively.\par
The coefficient of friction between the lamina and the wall is μ.\par
\textbf{(i)} Show that, when the lamina is in limiting equilibrium with the frictional force acting\par
upwards,\par
\[d sin(θ + φ) = (cos θ + μ sin θ)(a cos φ + b sin φ) . (∗)\]
\textbf{(ii)} How should (∗) be modified if the lamina is in limiting equilibrium with the frictional\par
force acting downwards?\par
\textbf{(iii)} Find a condition on d, in terms of a, b, tan θ and tan φ, which is necessary and sufficient\par
for the frictional force to act upwards. Show that this condition cannot be satisfied if\par
\[b(2 tan θ + tan φ) < a.\]
\bigskip\textbf{Question 10.} A particle is projected from a point O on horizontal ground with initial speed u and at an\par
angle of θ above the ground. The motion takes place in the x-y plane, where the x-axis is\par
horizontal, the y-axis is vertical and the origin is O. Obtain the Cartesian equation of the\par
\[particle's trajectory in terms of u, g and λ, where λ = tan θ.\]
Now consider the trajectories for different values of θ with u fixed. Show that for a given\par
value of x, the coordinate y can take all values up to a maximum value, Y , which you should\par
determine as a function of x, u and g.\par
Sketch a graph of Y against x and indicate on your graph the set of points that can be\par
reached by a particle projected from O with speed u.\par
Hence find the furthest distance from O that can be achieved by such a projectile.\par
\end{flushleft}
\subsection*{Page 6}
\begin{flushleft}
\bigskip\textbf{Question 11.} A small smooth ring R of mass m is free to slide on a fixed smooth horizontal rail. A light\par
inextensible string of length L is attached to one end, O, of the rail. The string passes\par
\[through the ring, and a particle P of mass km (where k > 0) is attached to its other end;\]
this part of the string hangs at an acute angle α to the vertical and it is given that α is\par
constant in the motion.\par
Let x be the distance between O and the ring. Taking the y-axis to be vertically upwards,\par
write down the Cartesian coordinates of P relative to O in terms of x, L and α.\par
\textbf{(i)} By considering the vertical component of the equation of motion of P , show that\par
\[kmx¨ cos α = T cos α - kmg ,\]
where T is the tension in the string. Obtain two similar equations relating to the\par
horizontal components of the equations of motion of P and R.\par
sin α\par
\textbf{(ii)} Show that = k, and deduce, by means of a sketch or otherwise, that\par
(1 - sin α)2\par
motion with α constant is possible for all values of k.\par
\textbf{(iii)} Show that x¨ = -g tan α .\par
\end{flushleft}
\subsection*{Page 7}
\begin{flushleft}
\bigskip\textbf{Section C: Probability and Statistics}\par
\bigskip\textbf{Question 12.} The lifetime of a fly (measured in hours) is given by the continuous random variable T\par
with probability density function f(t) and cumulative distribution function F(t). The hazard\par
\[function, h(t), is defined, for F(t) < 1, by\]
f(t)\par
\[h(t) = .\]
1 - F(t)\par
\textbf{(i)} Given that the fly lives to at least time t, show that the probability of its dying within\par
the following δt is approximately h(t) δt for small values of δt.\par
\textbf{(ii)} Find the hazard function in the case F(t) = t/a for 0 < t < a. Sketch f(t) and h(t) in\par
this case.\par
\textbf{(iii)} The random variable T is distributed on the interval t > a, where a > 0, and its hazard\par
function is t\par
-1.\par
Determine the probability density function for T .\par
\textbf{(iv)} Show that h(t) is constant for t > b and zero otherwise if and only if f(t) = ke\par
-k(t-b)\par
\[for t > b, where k is a positive constant.\]
\textbf{(v)} The random variable T is distributed on the interval t > 0 and its hazard function is\par
given by\par
λ\par
\[h(t) =\]
tλ-1\par
,\par
\[\int θλ \sum\]
where λ and θ are positive constants. Find the probability density function for T .\par
\end{flushleft}
\subsection*{Page 8}
\begin{flushleft}
\bigskip\textbf{Question 13.} A random number generator prints out a sequence of integers I , I , I , . . . . Each integer\par
1 2 3\par
is independently equally likely to be any one of 1, 2, . . . , n, where n is fixed. The random\par
variable X takes the value r, where I is the first integer which is a repeat of some earlier\par
r\par
integer.\par
\[Write down an expression for P(X = 4).\]
\textbf{(i)} Find an expression for P(X = r), where 2 \leq r \leq n + 1. Hence show that, for any\par
positive integer n,\par
1 1 2 1 2 3\par
\[+ 1 - + 1 - 1 - + \cdots = 1 .\]
\[n \int n \sum n \int n \sum \int n \sum n\]
\textbf{(ii)} Write down an expression for E(X). (You do not need to simplify it.)\par
\textbf{(iii)} Write down an expression for P(X \geq k).\par
\textbf{(iv)} Show that, for any discrete random variable Y taking the values 1, 2, . . . , N ,\par
N\par
\[E(Y ) = P(Y \geq k) .\]
\[)k=1\]
Hence show that, for any positive integer n,\par
12 1 22 1 2 32\par
\[1 - + 1 - 1 - + 1 - 1 - 1 - + \cdots = 0.\]
\[\int n \sum \int n \sum \int n \sum \int n \sum \int n \sum \int n \sum\]
\end{flushleft}
\end{document}
